\section{USAMO 1989}
        \begin{enumerate}
        	\item (\textit{USAMO 1989}) For each positive integer $n$, let


 <div style="text-align:center;">
$S_n = 1 + \frac 12 + \frac 13 + \cdots + \frac 1n$


 $T_n = S_1 + S_2 + S_3 + \cdots + S_n$


 $U_n = \frac{T_1}{2} + \frac{T_2}{3} + \frac{T_3}{4} + \cdots + \frac{T_n}{n+1}$.


 
</div>Find, with proof, integers $0 < a,\ b,\ c,\ d < 1000000$ such that $T_{1988} = a S_{1989} - b$ and $U_{1988} = c S_{1989} - d$.


 



\textbf{Solution}

We note that for all integers $n \ge 2$,
 \begin{align*} T_{n-1} &= 1 + \left(1 + \frac 12\right) + \left(1 + \frac 12 + \frac 13\right) + \ldots + \left(1 + \frac 12 + \frac 13 + \ldots + \frac 1{n-1}\right) \\  &= \sum_{i=1}^{n-1} \left(\frac {n-i}i\right) = n\left(\sum_{i=1}^{n-1} \frac{1}{i}\right) - (n-1) = n\left(\sum_{i=1}^{n} \frac{1}{i}\right) - n \\ &= n \cdot S_{n} - n .  \end{align*}


 It then follows that
 \begin{align*} U_{n-1} &= \sum_{i=2}^{n} \frac{T_{i-1}}{i} = \sum_{i=2}^{n}\ (S_{i} - 1) = T_{n-1} + S_n - (n-1) - S_1 \\ &= \left(nS_n - n\right) + S_n - n = (n + 1)S_n - 2n . \end{align*}


 If we let $n=1989$, we see that $(a,b,c,d) = (1989,1989,1990, 2\cdot 1989)$ is a suitable solution.  $\blacksquare$

 Notice that it is also possible to use induction to prove the equations relating $T_n$ and $U_n$ with $S_n$.

 
<i>Alternate solutions are always welcome. If you have a different, elegant solution to this problem, please \href{https://artofproblemsolving.com/wiki/index.php?title=1989\_USAMO\_Problems/Problem\_1&action=edit}{add it} to this page.</i>

 


	\item (\textit{USAMO 1989}) The 20 members of a local tennis club have scheduled exactly 14 two-person games among themselves, with each member playing in at least one game. Prove that within this schedule there must be a set of 6 games with 12 distinct players.


 



\textbf{Solution 1}

Consider a graph with $20$ vertices and $14$ edges. The sum of the degrees of the vertices is $28$; by the \href{https://artofproblemsolving.com/wiki/index.php?title=Pigeonhole\_Principle}{Pigeonhole Principle} at least $12$ vertices have degrees of $1$ and at most $8$ vertices have degrees greater than $1$. If we keep deleting edges of vertices with degree greater than $1$ (a maximum of $8$ such edges), then we are left with at least $6$ edges, and all of the vertices have degree either $0$ or $1$. These $6$ edges represent the $6$ games with $12$ distinct players.

 



\textbf{Solution 2}

Let a slot be a place we can put a member in a game, so there are two slots per game, and 28 slots total.  We begin by filling exactly 20 slots each with a distinct member since each member must play at least one game.  Let there be $m$ games with both slots filled and $n$ games with only one slot filled, so $2m+n=20$.  Since there are only 14 games, $m+n \leq 14 \Longrightarrow 2m+n \leq 14+m \Longleftrightarrow 20 \leq 14+m \Longrightarrow m \geq 6$, so there must be at least 6 games with two distinct members each, and we must have our desired set of 6 games.

 



\textbf{Solution 3}

Assume the contrary. 

 Consider the largest set of disjoint edges $E$. By assumption it has less than $6$ edges, i.e. maximum $10$ vertices. Call it a vertex set $V$. 

 $10$ vertices remain outside $V$ and each has to be attached to at least one edge. Now, if any two vertices outside $V$ are connected by, say, edge $e$, we could have included $e$ in $E$ and gotten a larger disjoint set, so - a contradiction. Therefore the only option would be that all vertices outside $V$ are connected each by one edge to some vertices inside $V$. That would take $10$ edges, but $E$ already includes $5$ - again a contradiction. 

 All possibilities yield a contradiction, so our assumption can not be correct.

 (Cases when largest set $E$ is smaller than $6$ are equivalent and weaker)

 



\textbf{Solution 4}

Assume, for the sake of contradiction, that there is no set of 6 games with 12 distinct players. Then the largest possible collection of games with all distinct players must consist of at most 5 games, involving at most 10 players.So suppose we have 5 such games with 10 distinct players. This leaves 10 other players who must still appear in the schedule. Since there are a total of 14 games, there are $14-5=9$ games remaining to cover these 10 players.

 Each game has 2 players, so the 9 games account for 18 ``player slots.\textit{ If every one of the 10 unused players is to appear at least once, then by the pigeonhole principle, some game among these 9 must contain two of the 10 unused players.}But then that game, together with the original 5 disjoint games, would give 6 games with 12 distinct players. This contradicts our assumption that no such collection exists.Therefore, there must in fact exist a set of 6 games with 12 distinct players.

 


	\item (\textit{USAMO 1989}) Let $P(z)= z^n + c_1 z^{n-1} + c_2 z^{n-2} + \cdots + c_n$ be a polynomial in the complex variable $z$, with real coefficients $c_k$. Suppose that $|P(i)| < 1$. Prove that there exist real numbers $a$ and $b$ such that $P(a + bi) = 0$ and $(a^2 + b^2 + 1)^2 < 4 b^2 + 1$.


 



\textbf{Solution}

Let $z_1, \dotsc, z_n$ be the (not necessarily distinct) roots of $P$, so that
 \[P(z) = \prod_{j=1}^n (z- z_j) .\]Since all the coefficients of $P$ are real, it follows that if $w$ is a root of $P$, then $P( \overline{w}) = \overline{ P(w)} = 0$, so $\overline{w}$, the \href{https://artofproblemsolving.com/wiki/index.php?title=Complex\_conjugate}{complex conjugate} of $w$, is also a root of $P$.

 Since
 \[\lvert i- z_1 \rvert \cdot \lvert i - z_2 \rvert  \dotsm \lvert i - z_n \rvert = \lvert P(i) \rvert < 1,\]it follows that for some (not necessarily distinct) conjugates $z_i$ and $z_j$,
 \[\lvert z_i-i \rvert \cdot \lvert z_j-i \rvert < 1.\]Let $z_i = a+bi$ and $z_j = a-bi$, for real $a,b$.  We note that
 \[(a+b+1)^2 - (a+b-1)^2 = 4a + 4b .\]Thus
 \begin{align*} (a^2+b^2+1)^2 &= (a^2+b^2-1)^2 + 4a^2 + 4b^2 = \lvert a^2 + b^2 - 1 - 2ai \rvert ^2 + 4b^2 \\ &= \lvert (a-i)^2 - (bi)^2 \rvert^2 + 4b^2 \\ &= \bigl( \lvert a+bi - i \rvert \cdot \lvert a-bi -i \rvert \bigr)^2 + 4b^2 \\ &= \bigl( \lvert z_i - i \rvert \cdot \lvert z_j - i \rvert \bigr)^2 + 4b^2 < 1+4b^2. \end{align*}
Since $P(a+bi) = P(z_i) = 0$, these real numbers $a,b$ satisfy the problem's conditions.  $\blacksquare$

 


	\item (\textit{USAMO 1989}) Let $ABC$ be an acute-angled triangle whose side lengths satisfy the inequalities $AB < AC < BC$. If point $I$ is the center of the inscribed circle of triangle $ABC$ and point $O$ is the center of the circumscribed circle, prove that line $IO$ intersects segments $AB$ and $BC$.


 



\textbf{Solution 1}

Proceed with barycentrics. First, note that the desired conclusion is equivalent to showing that $X=\overleftrightarrow{IO} \cap \overleftrightarrow{AC}$ is not on segment $AC$. 

 Let $X=(a:0:t)$, so 
 \[0 = \begin{vmatrix} a & 0 & t \\ a & b & c \\ a^2S_A & b^2S_B & c^2S_C \end{vmatrix} = abc \begin{vmatrix} 1 & 0 & \frac{t}{c} \\ 1 & 1 & 1 \\ aS_A & bS_B & cS_C \end{vmatrix} \implies t=\frac{c(cS_C-bS_B)}{aS_A-bS_B}\]But, 
 \[t = \frac{c(cS_C-bS_B)}{aS_A-bS_B}=\frac{c(c(a^2+b^2-c^2)-b(a^2+c^2-b^2))}{a(b^2+c^2-a^2)-b(a^2+c^2-b^2)}=\frac{c(c-b)(a^2+b^2+c^2)}{(a-b)(a^2+b^2+c^2)}=\frac{c(c-b)}{a-b} < 0\]as needed.

 -sillybone

 



\textbf{Solution 2}

Consider the lines that pass through the circumcenter $O$. Extend $AO$, $BO$, $CO$ to $D$,$E$,$F$ on $a$,$b$,$c$, respectively.

 We notice that $IO$ passes through sides $a$ and $c$ if and only if $I$ belongs to either regions $AOF$ or $COD$.

 


 Since $AO = BO = CO = R$, we let $\alpha = \angle OAC = \angle OCA$, $\beta = \angle BAO = \angle ABO$, $\gamma = \angle BCO = \angle CBO$.

 We have $c < b < a\implies C < B < A\implies$ $\gamma+\alpha <\beta+\gamma <\alpha+\beta\implies\gamma <\alpha <\beta$

 Since $IA$ divides angle $A$ into two equal parts, it must be in the region marked by the $\beta$ of angle $A$, so $I$ is in $ABD$.

 Similarly, $I$ is in $ACF$ and $ABE$. Thus, $I$ is in their intersection, $AOF$. From above, we have $IO$ passes through $a$ and $c$. $\blacksquare$

 


	\item (\textit{USAMO 1989}) Let $u$ and $v$ be real numbers such that


 $(u + u^2 + u^3 + \cdots + u^8) + 10u^9 = (v + v^2 + v^3 + \cdots + v^{10}) + 10v^{11} = 8.$


 Determine, with proof, which of the two numbers, $u$ or $v$, is larger.


 



\textbf{Solution}

The answer is $v$.

 We define real functions $U$ and $V$ as follows:
 \begin{align*} U(x) &= (x+x^2 + \dotsb + x^8) + 10x^9 = \frac{x^{10}-x}{x-1} + 9x^9 \\ V(x) &= (x+x^2 + \dotsb + x^{10}) + 10x^{11} = \frac{x^{12}-x}{x-1} + 9x^{11} . \end{align*}
We wish to show that if $U(u)=V(v)=8$, then $u <v$.

 We first note that when $x \le 0$, $x^{12}-x \ge 0$, $x-1 < 0$, and $9x^9 \le 0$, so
 \[U(x) = \frac{x^{10}-x}{x-1} + 9x^9 \le 0 < 8 .\]Similarly, $V(x) \le 0 < 8$.

 We also note that if $x \ge 9/10$, then
 \begin{align*} U(x) &= \frac{x-x^{10}}{1-x} + 9x^9 \ge \frac{9/10 - 9^9/10^9}{1/10} + 9 \cdot \frac{9^{9}}{10^9} \\ &= 9 - 10 \cdot \frac{9^9}{10^9} + 9 \cdot \frac{9^9}{10^9} = 9 - \frac{9^9}{10^9} > 8. \end{align*}
Similarly $V(x) > 8$.  It then follows that $u, v \in (0,9/10)$.

 Now, for all $x \in (0,9/10)$,
 \begin{align*} V(x) &= U(x) + V(x)-U(x) = U(x) + 10x^{11}+x^{10} -9x^9 \\ &= U(x) + x^9 (10x -9) (x+1) < U(x) . \end{align*}
Since $V$ and $U$ are both strictly increasing functions over the nonnegative reals, it then follows that
 \[V(u) < U(u) = 8 = V(v),\]so $u<v$, as desired.  $\blacksquare$

 


\end{enumerate}